% Floods are the most frequent natural disaster and have impacted every U.S. state and nearly every country. The
% connection between floods and climate change comes down to a few ways that climate change is impacting water. First,
% higher temperatures lead to increased levels of evaporation, creating denser clouds that hold more water. This
% eventually leads to heavier precipitation that can cause flooding. Second, more frequent and intense storms such as
% hurricanes can lead to floods. Finally, higher sea levels due to melting glaciers can also prompt coastal flooding.

% 1. Aumento de temperaturas globais;
% 2. Ação humana no clima;
% 3. Aumento de heat waves -> Dunn;
% 4. Aumento de precipitação;
% 5. Relação disso com aumento de desastres naturais em geral;
% 6. Aumento de Enchentes (precipitação/precipitação acumulada);
% 7. Sul Global mais afetado proporcionalmente

% Chapter 2.3.1.1.3 (Temperatures during the instrumental period - surface) of the AR6 WGI

As últimas décadas foram marcadas por quebras constantes de recordes climáticos, especialmente em máximas de
temperatura, evidenciando o impacto significativo das mudanças climáticas em escala global (EXPANDIR). De acordo com o
sexto - e mais recente - relatório do Painel Intergovernamental sobre Mudanças Climáticas (IPCC), estabelecido pelas
Nações Unidas e amplamente reconhecido como a principal referência nas análises do conhecimento científico relacionado
às mudanças climáticas, a superfície terrestre tem experimentado um aumento expressivo de temperatura em comparação com
períodos históricos. Considerando o intervalo de 1850 a 1900 \footnote{Utilizado como base pré-industrial. É utilizado
como período de referencia por ser o primeiro com boas observações de temperaturas terrestres registradas}, a
temperatura média da superfície terrestre\footnote{Dados referentes à medida do GMST - Temperatura Média Global da
Superfície - que representa uma média ponderada entre as temperaturas das massas terrestres e das massas oceânicas.} no
período de 2001 a 2020 foi aproximadamente 1ºC mais elevada. Além disso, a tendência de aumento da temperatura somente
nas últimas quatro décadas (1980 a 2020) é de 0,76ºC, evidenciando uma aceleração no aquecimento global \cite{gulev}.

Esse aumento, aparentemente pequeno, não reflete a real magnitude das modificações no sistema terrestre como um todo. É
importante notar que um aumento de 1ºC representa uma média global; portanto, houve locais em que o aquecimento foi
significativamente mais pronunciado, como nos polos. Nessas regiões, qualquer aquecimento pode levar ao derretimento de
geleiras e, consequentemente, à diminuição do albedo, criando um ciclo de retroalimentação que acelera ainda mais o
aquecimento (EXPANDIR). Além disso, é relevante destacar que um aquecimento de 1ºC em toda a superfície terrestre é
extremamente significativo devido à massa do planeta. Para que um aumento dessa magnitude ocorra em um corpo com a massa
da Terra, a quantidade de energia térmica\footnote{De forma bastante simplificada, essa energia pode ser vista pela
fórmula $ Q = mc\Delta T $, onde $m$ é a massa, $c$ é o calor específico, e $\Delta T$ a variação de temperatura.}
necessária é imensa.

% Acordo de Paris?

% Chapter 3.3 (Human Influence on the Atmosphere and Surface) of the AR6 WGI;
É indubitável que as atividades humanas tem sido um grande contribuinte para o Aquecimento Global.

% A ação humana no frágil clima terrestre é indubutável e sentida constanemente de forma empirca. tem implicações
% profundas nos sistemas climáticos globais, influenciando eventos extremos como ondas de calor, secas prolongadas, e
% padrões de precipitação irregular.

% Este dado reflete a intensificação das atividades humanas que contribuem para o efeito estufa, como a queima de
% combustíveis fósseis, desmatamento e emissões industriais. Esses números, embora alarmantes, reforçam a necessidade
% urgente de ações concretas e coordenadas em escala global para mitigar os impactos climáticos e promover um futuro
% mais sustentável 

% Grande parte da mudança na América Latina é devido à LULUCF, i.e., uso da terra -> Plantações, desmatamento e outros;

% Ou seja, a América Latina sempre foi um grande contribuinte para o aquecimento global mas não como grande consumidor
% das suas vantagens, mas sim como fornecedor de suas demandas.